\chapter{Electronic dance music: waking the body first}
\label{ch:genre-edm-deep}

\classquote{乐者，音之所由生也；其本在人心之感于物也。}{《礼记·乐记》}{Music is born from sound, but its root is a heart moved by the world --- dance music makes the body nod before the mind understands.}

\section{What EDM is: a machine for building anticipation}
\label{sec:edm-what}

\term{electronic dance music}{electronic dance music} is not ``music made with a computer''. It has no fixed lineup and is not defined by any single instrument. It defines itself by an agreement about how the listener's \emph{body} will take part: a four-on-the-floor kick lays down a predictable floor, offbeat hats lift the body up, and the whole song keeps answering one question --- ``will the next beat bring something bigger?''

You probably know the feeling. A song plays in your headphones, the intro is ordinary, then every sound is pulled away until only a rising noise remains. Two seconds later the low end crashes back and something in you sinks. That is not an accident; it is design. EDM turns ``anticipation and payoff'' into a repeatable structure, and the \term{drop}{drop} is the moment of payoff. Its charm is not how cool the sounds are, but that it knows exactly when you will hold your breath.

This makes EDM especially inviting for young people to make on their own. You need no acoustic instrument and no rehearsal room --- just a computer, a \term{digital audio workstation}{digital audio workstation}, and a few loops to build something a friend will nod along to. The entry barrier is low, but the ceiling is high: the hard part is never ``putting a sound in'', it is ``knowing when to leave one out''.

Break the listening experience down and four things hold it up:

\begin{longtable}{p{0.16\textwidth}p{0.32\textwidth}p{0.38\textwidth}}
\toprule
Handle & What it feels like & Question to write against \\
\midrule
\endfirsthead
\toprule
Handle & What it feels like & Question to write against \\
\midrule
\endhead
Pulse & The four-on-the-floor kick is a heartbeat; the body moves before the mind & Which beat is the floor? Is it there every time? \\
Breath & Sounds are pulled away together and return together, like a tide & Is there a subtraction here, so the return hits harder? \\
Loop & A short motif returns, changing only a little each time & What is the smallest memorable unit, and what changed on repeat? \\
Impact & Low end and main body land together in a moment of release & Was the anticipation earned? Is the payoff bigger than the promise? \\
\bottomrule
\end{longtable}

What EDM fears is not ``loud'' but ``always full''. If every beat sounds and every layer moves, the listener has nothing left to anticipate. Professionalism often comes from restraint: knowing which layer should step back, which beat should be empty, and which return should keep one trick in reserve.

\section{From Chicago warehouses to the main stage}
\label{sec:edm-history}

EDM did not appear overnight. It inherited the dance-floor logic of \term{disco}{disco}, learned skeletal arrangement from \term{house}{house}, borrowed the idea of timbre-as-hook from \term{synth-pop}{synth-pop}, and grew new branches with every change in equipment and distribution. The table below follows only the main line where ``the technology or the scene changed, so the arrangement changed too''; it does not try to list every branch.

\begin{longtable}{p{0.16\textwidth}p{0.34\textwidth}p{0.42\textwidth}}
\toprule
Period & Key change & Effect on arrangement \\
\midrule
\endfirsthead
\toprule
Period & Key change & Effect on arrangement \\
\midrule
\endhead
1970s & New York and Chicago clubs, disco and drum machines & Four-on-the-floor becomes a body coordinate; DJs start editing records like arrangers \\
1980s & Chicago house, Detroit techno, widely available synths and samplers & Sound design becomes part of arranging; \term{frequency-band division}{frequency-band division} replaces ``sections of a band'' \\
1990s & European dance music, trance and large parties & Sections grow longer; \term{build-up}{build-up} and drop become the structural core \\
2000s & French house, digital production and podcasts & Filters become a lead character; a track can be finished at home and tested in a club \\
2010s & Festival EDM, the stage economy and streaming & Drops get brighter and more direct; even a fifteen-second clip must work \\
2020s & Short video, recommendation algorithms and genre blending & Intros shrink; vintage synth tones and modern low end share one song \\
\bottomrule
\end{longtable}

(The historical outline follows Wikipedia's ``Electronic dance music'' and ``House music'' articles; see further reading for named artists and years.)

What is interesting is how often EDM looks back. Many of the biggest hits of the 2020s deliberately use 1980s synth tones with today's heavier low end. Genres do not improve in a straight line --- they keep reinterpreting old sounds as they move forward.

\section{Arranging: divide by frequency, not by instrument section}
\label{sec:edm-craft}

A traditional arrangement asks ``who plays what''; EDM first asks ``which frequency band owns this moment''. The kick owns the lowest transients, the \term{sub-layer}{sub-layer} owns sustained low frequencies, the bass writes in the lower mids, the pad fills the mids, and the lead and pluck carry melody and air. Once every layer signs a ``frequency contract'', mixing becomes half the work.

\subsection{Typical setup and four functional layers}
\label{sec:edm-instrumentation}

\begin{longtable}{p{0.16\textwidth}p{0.26\textwidth}p{0.26\textwidth}p{0.22\textwidth}}
\toprule
Layer & Typical tracks & Main job & Common mistake \\
\midrule
\endfirsthead
\toprule
Layer & Typical tracks & Main job & Common mistake \\
\midrule
\endhead
Impact & Kick, clap, snare & Predictable floor and accents & Every layer fights for the transient, so the kick disappears \\
Low end & Sub, bass & Sustained low frequencies and roots & Sub and bass share one register and blur into mud \\
Harmony & Pad, chord, arp & Bed and light/dark contrast & Pad fills the whole spectrum, leaving no room for the lead \\
Melody and air & Lead, pluck, FX & Memorable hooks and section edges & They play from start to finish, so the drop is no longer fresh \\
\bottomrule
\end{longtable}

\subsection{Groove: four-on-the-floor, offbeats and sidechain}
\label{sec:edm-groove}

\term{Four-on-the-floor}{four-on-the-floor} means the kick hits on every beat; offbeat open hats lift the body; a clap or snare on beats 2 and 4 gives something to shout along with. These three layers look simple but form the most stable coordinate --- melody and timbre can fly anywhere as long as the floor does not move.

\term{Sidechain}{sidechain} is EDM's own way of breathing: when the kick hits, the pad and bass duck slightly, then return. It feels like the whole band is \term{pumping}{pumping}. This is not an effect for its own sake; it is an arranging device that keeps the low-end layer and the impact layer from fighting. In the app, bind a pad track's volume automation to the kick pattern, or use a compressor's sidechain input.

\subsection{Harmony and melody: polishing one loop}
\label{sec:edm-harmony}

EDM often uses a single four- or eight-bar loop --- vi--IV--I--V or i--VI--III--VII, for instance. The harmony does not aim to be ``rich''; it aims to be stable enough to forget. The freed attention goes to timbre and rhythm: the same chords first on a plain pad, then with a pluck, then with the filter's \term{cutoff frequency}{cutoff frequency} slowly raised --- the listener hears ``the same song growing up''.

A lead need not be complex. A two-bar motif, repeated, transposed and re-rhythmed, is enough to carry a section. What matters is that the motif is short enough to remember and leaves one place where something can change.

\subsection{Timbre and automation: making sound move}
\label{sec:edm-automation}

In the app, automation is not a finishing touch; it is part of arranging:

\begin{itemize}
\item Raise the pad's filter cutoff slowly from the start of a section to its brightest just before the drop.
\item Automate the volume of a noise or riser upward through the build, then pull it back on the first beat of the drop.
\item Pan an arpeggio from one side to the other to add spatial motion to a static loop.
\end{itemize}

\subsection{Section structure: subtract first, then add}
\label{sec:edm-structure}

A common frame is Intro (8 bars) $\to$ Groove (16) $\to$ Build (8) $\to$ Drop (16) $\to$ Break (8) $\to$ Build (8) $\to$ Drop (16) $\to$ Outro, with eight bars per block and no chord changes inside a block.

The key idea: a \term{build-up}{build-up} is not ``add everything''; it is ``take the floor away, then give it back''. The impact of a drop comes from those two empty seconds before it. If you push every layer to maximum during the build, the drop has nowhere left to go.

\section{Four tracks, taken apart}
\label{sec:edm-cases}

The four below are two classics and two recent favourites, chosen with younger listeners in mind. The analysis uses only arrangement and orchestration language: no lyrics and no full melody transcription. Names and dates are sourced in the further-reading list.

\subsection{Daft Punk, ``One More Time'': the filter is the lead character}
\label{sec:edm-case-daft}

The 2000 track ``One More Time'' by the French duo Daft Punk, from the album \emph{Discovery}, is a textbook case of French house. Its most famous move is turning a filter sweep into an arrangement event: a sound that seems buried behind a wall is slowly opened until it floats forward. Here, changing timbre \emph{is} the section --- no new instrument is needed.

The orchestration is restrained: drums, bass, chords and a treated vocal, each with a clear band. When you listen, count just one thing --- how many bars does the filter take to go from muffled to bright? Copying that number into your own song is more useful than copying the preset.

\subsection{Avicii, ``Levels'': one pluck holds a whole song}
\label{sec:edm-case-avicii}

Swedish producer Avicii's 2011 single ``Levels'' proves that an EDM hook can be one line. A recurring pluck motif over a clean four-on-the-floor is enough to make the whole track memorable. It works not through complexity but through \term{variation within repetition}{variation within repetition} --- each return changes the timbre or rhythmic density slightly, so the ear never tires.

(Source: Avicii, ``Levels'', 2011 single; see the further-reading links for the official release.)

\subsection{Martin Garrix, ``Animals'': the drop as a product}
\label{sec:edm-case-garrix}

Dutch producer Martin Garrix's 2013 single ``Animals'' is the most direct example of the big-room era: a long build, a short hook, and a huge, simple drop. Its arrangement hides almost nothing --- low end and main body land together, and the only silence is just before impact.

What is worth learning is its economy: the track does not stack many layers; it saves two or three sounds for the most important moments. A common mistake for young producers is wanting something new in every bar. This song shows the opposite kind of courage.

\subsection{The Weeknd, ``Blinding Lights'': vintage tone, modern low end}
\label{sec:edm-case-weeknd}

The Weeknd's 2019 track ``Blinding Lights'', from \emph{After Hours}, joins 1980s synth-pop tones to contemporary dance low end and became a recent global hit. Its drums are thoroughly modern while its timbres deliberately carry the texture of old machines; two senses of time overlap in one song, and that feels fresh.

It is a reminder that genres evolve by re-placing old sounds in new contexts, not by erasing them. When you try something similar in the app, write the harmony layer with vintage tones first, then write the kick and sub with modern sounds --- let each layer stand on its own, then listen to see whether they can coexist.

\section{The app workflow: from blank session to first drop}
\label{sec:edm-yeban}

Build in the app in this order, solving one problem at a time:

\begin{enumerate}
\item Create tracks: Kick (audio or MIDI), Sub (MIDI), Bass (MIDI), Pad (MIDI), Lead (MIDI), Hats (MIDI), and a Master bus. Seven tracks are enough; clear the workspace first.
\item Write the floor: a kick on every beat, offbeat eighth hats, clap on 2 and 4. Loop only these three layers and check that your body nods.
\item Write the low end: long root notes on Sub, short eighths on Bass. Keep them out of the same register; separate them by \term{frequency-band division}{frequency-band division}.
\item Write harmony and motif: a four-bar pad loop and a two-bar lead motif. Write only these two, with no tone shaping yet.
\item Shape the rise: automate the pad's cutoff from dark to bright, replace the drums with an accelerating snare or noise in the build, and leave two bars of silence before the drop.
\item Assemble the form: eight bars per block --- Intro--Groove--Build--Drop--Break--Build--Drop--Outro. For the first pass, stay under 64 bars.
\item Check: listen for muddy low end, for the pad fighting the lead's band, and for a drop that is genuinely bigger than what came before. If something is wrong, remove a layer before adding one.
\end{enumerate}

\section{Practice: three routes, different trade-offs}
\label{sec:edm-practice}

\begin{exercise}
\textbf{One loop, three drops.} With the same vi--IV--I--V loop and the same kick, make three drops: A pulls every layer out for two beats then brings them all back; B brings back only low end and body, letting the melody wait one bar; C uses a filter sweep instead of a volume impact. Compare ``impact, audible detail, and editability''. A is the most direct but can be brutal; B has more structure but watch that the groove does not loosen; C is the most subtle but depends on sound-design skill.
\end{exercise}

\begin{exercise}
\textbf{A two-bar motif, three ways to repeat.} Write a two-bar lead motif and make three versions: A repeats it exactly; B changes the rhythm of the last note each time; C transposes it or changes its timbre. Ask two people to listen without seeing the screen, then ask which version they remember without getting tired of it. A is stable but hypnotic; B changes little yet sounds most human; C carries more information and may pull attention from the low end.
\end{exercise}

\begin{exercise}
\textbf{Three ways to pump.} Give the same pad a pumping feel: A draw volume automation against the kick in the app; B use a compressor's sidechain input; C manually shorten the pad at each kick. Compare ``precision, speed, and portability''. A is visible but slow to draw; B is the regular studio approach; C is crude but works in any software. Try all three before deciding which one your song needs.
\end{exercise}

\section{Further reading}
\label{sec:edm-reading}

\begin{sloppypar}
\begin{itemize}
\item \textbf{Wikipedia, ``House music''}: \url{https://en.wikipedia.org/wiki/House_music}. Trace Chicago house into EDM, filling in the details of this chapter's historical line.
\item \textbf{Wikipedia, ``Electronic dance music''}: \url{https://en.wikipedia.org/wiki/Electronic_dance_music}. A chronological pass from house, techno and trance to EDM, for checking the historical line of this chapter.
\item \textbf{Daft Punk, ``One More Time'' (Official Video)}: \url{https://www.youtube.com/watch?v=FGBhQbmPwH8}. Hear how a filter sweep takes the place of ``adding another instrument'' as a section event.
\item \textbf{Avicii, ``Levels'' (Official Video)}: \url{https://www.youtube.com/watch?v=_ovdm2yX4MA}. Hear how one pluck motif carries a whole song through variation within repetition.
\item \textbf{Martin Garrix, ``Animals'' (Official Video)}: \url{https://www.youtube.com/watch?v=g9WH3n36WVE}. See how the big-room build and drop achieve their effect with extreme economy.
\item \textbf{The Weeknd, ``Blinding Lights'' (Official Video)}: \url{https://www.youtube.com/watch?v=4NRXx6U8ABQ}. Compare the mix of vintage synth tones with contemporary low end.
\end{itemize}
\end{sloppypar}

Links above were checked against official pages and search results on 27 November 2026 for titles and authorship; if a platform's regional restriction prevents playback, search for the item on the artist's official account.

\section{Takeaways}
\label{sec:edm-takeaways}

\begin{itemize}
\item EDM is not ``music made with a computer'' but a machine for building anticipation: four-on-the-floor gives the floor, sidechain gives the breath, and the drop gives the payoff.
\item Genres follow equipment and venues --- from clubs and drum machines to streaming and short video, each era changes arranging priorities.
\item Write frequency contracts before instruments; each layer does one job, and track count is not a measure of density.
\item Harmony seeks stability, melody seeks memorability; change goes into timbre, rhythmic density and automation.
\item A build subtracts first; the impact of a drop comes from the silence before it.
\item The app workflow runs ``floor --- low end --- harmony and motif --- automation --- structure --- check''.
\item Practice compares trade-offs; an arrangement is usable when you can change it, delete it, and explain it.
\end{itemize}

\begin{helpmapping}
\begin{longtable}{p{0.45\textwidth}p{0.45\textwidth}}
\toprule
Section & Help anchor \\
\midrule
What EDM is & \texttt{arrangement-genre-edm-what} \\
From Chicago to the main stage & \texttt{arrangement-genre-edm-history} \\
Arranging & \texttt{arrangement-genre-edm-craft} \\
Four tracks, taken apart & \texttt{arrangement-genre-edm-cases} \\
The app workflow & \texttt{arrangement-genre-edm-yeban} \\
Practice & \texttt{arrangement-genre-edm-practice} \\
Further reading & \texttt{arrangement-genre-edm-reading} \\
\bottomrule
\end{longtable}
\end{helpmapping}

\begin{figure}[htbp]\centering
\begin{tikzpicture}[font=\scriptsize,x=1cm,y=1cm,>=Stealth]
\yebangenreddraw{EDM: shape energy into a waveform}{20--60 Hz sub}{120--2 kHz lead}{8 kHz+ air}
\end{tikzpicture}
\caption{EDM frequency split, drum grid and section-energy curve}
\end{figure}
